% from ega1-7-fr.tex
\oldpage[I]{158}
\begin{proposition}[7.1.11]
\label{I.7.1.11-fr}
Soient $X$, $Y$ deux $S$-préschémas~; on suppose $X$ irréductible, de point
générique $x$, et $Y$ de type fini sur $S$. Deux $S$-applications rationnelles
de $X$ dans $Y$, auxquelles correspond le même $S$-morphisme
$\operatorname{Spec}(\mathscr{O}_x)\to Y$ sont identiques. Si on suppose de
plus $S$ localement noethérien, tout $S$-morphisme de
$\operatorname{Spec}(\mathscr{O}_x)$ dans $Y$ correspond à une $S$-application
rationnelle (et une seule) de $X$ dans $Y$.
\end{proposition}

Compte tenu de ce que tout ouvert non vide de $X$ est partout dense, cela
résulte aussitôt de (\hyperref[I.6.5.1-fr]{6.5.1}).

\begin{corollary}[7.1.12]
\label{I.7.1.12-fr}
On suppose $S$ localement noethérien, et les autres hypothèses de
(\hyperref[I.7.1.11-fr]{7.1.11}) satisfaites. Les $S$-applications rationnelles
de $X$ dans $Y$ s'identifient alors aux points du $S$-préschéma $Y$, à valeurs
dans le $S$-préschéma $\operatorname{Spec}(\mathscr{O}_x)$.
\end{corollary}

Cela n'est autre que (\hyperref[I.7.1.11-fr]{7.1.11}), avec la terminologie
introduite dans (\hyperref[I.3.4.1-fr]{3.4.1}).

\begin{corollary}[7.1.13]
\label{I.7.1.13-fr}
On suppose remplies les conditions de (\hyperref[I.7.1.12-fr]{7.1.12}). Soit
$s$ l'image de $x$ dans $S$. La donnée d'une $S$-application rationnelle de
$X$ dans $Y$ équivaut à la donnée d'un point $y$ de $Y$ au-dessus de $s$, et
d'un $\mathscr{O}_s$-homomorphisme local
$\mathscr{O}_y\to\mathscr{O}_x=R(X)$.
\end{corollary}

Cela résulte de (\hyperref[I.7.1.11-fr]{7.1.11}) et de
(\hyperref[I.2.4.4-fr]{2.4.4}).

En particulier~:

\begin{corollary}[7.1.14]
\label{I.7.1.14-fr}
Sous les conditions de (\hyperref[I.7.1.12-fr]{7.1.12}), les $S$-applications
rationnelles de $X$ dans $Y$ ne dépendent (pour $Y$ donné) que du
$S$-préschéma $\operatorname{Spec}(\mathscr{O}_x)$ et en particulier restent
les mêmes lorsqu'on remplace $X$ par
$\operatorname{Spec}(\mathscr{O}_z)$, pour tout $z\in X$.
\end{corollary}

En effet, comme $z\in\overline{\{x\}}$, $x$ est le point générique de
$Z=\operatorname{Spec}(\mathscr{O}_z)$, et
$\mathscr{O}_{X,x}=\mathscr{O}_{Z,x}$.

Lorsque $X$ est intègre, $R(X)=\mathscr{O}_x=k(x)$ est un corps
(\hyperref[I.7.1.5-fr]{7.1.5})~; les corollaires précédents se spécialisent
alors en~:

\begin{corollary}[7.1.15]
\label{I.7.1.15-fr}
On suppose vérifiées les conditions de (\hyperref[I.7.1.12-fr]{7.1.12}) et en
outre que $X$ est intègre. Soit $s$ l'image de $x$ dans $S$. Alors les
$S$-applications rationnelles de $X$ dans $Y$ s'identifient aux points
géométriques de $Y\otimes_S k(s)$ à valeurs dans l'extension $R(X)$ de $k(s)$,
autrement dit chacune d'elles équivaut à la donnée d'un point $y\in Y$
au-dessus de $s$ et d'un $k(s)$-monomorphisme de $k(y)$ dans
$k(x)=R(X)$.
\end{corollary}

Les points de $Y$ au-dessus de $s$ s'identifient en effet à ceux de
$Y\otimes_S k(s)$ (\hyperref[I.3.6.3-fr]{3.6.3}) et les
$\mathscr{O}_s$-homomorphismes locaux $\mathscr{O}_y\to R(X)$ aux
$k(s)$-monomorphismes $k(y)\to R(X)$.

Plus particulièrement~:

\begin{corollary}[7.1.16]
\label{I.7.1.16-fr}
Soient $k$ un corps, $X$, $Y$ deux préschémas algébriques
(\hyperref[I.6.4.1-fr]{6.4.1}) sur $k$~; on suppose en outre $X$ intègre. Alors
les $k$-applications rationnelles de $X$ dans $Y$ s'identifient aux points
géométriques de $Y$ à valeurs dans l'extension $R(X)$ de $k$
(\hyperref[I.3.4.4-fr]{3.4.4}).
\end{corollary}

\subsection{Domaine de définition d'une application rationnelle.}
\label{subsection:I.7.2-fr}

\begin{env}[7.2.1]
\label{I.7.2.1-fr}
Soient $X$, $Y$ deux préschémas, $f$ une application rationnelle de $X$ dans
$Y$. On dit que $f$ est \emph{définie en un point $x\in X$} s'il existe un
ensemble ouvert partout dense $U$ contenant $x$ et un morphisme $U\to Y$
appartenant à la classe d'équivalence $f$. L'ensemble des points $x\in X$ où
$f$ est définie est appelé le \emph{domaine de définition de $f$}~; il est
clair que c'est un ouvert partout dense dans $X$.
\end{env}

